\documentclass[12pt]{article}
\usepackage{graphicx,xcolor}
\usepackage[hidelinks]{hyperref}
\usepackage[margin=0.8in]{geometry}
\usepackage{listings}
\lstset{
    basicstyle=\ttfamily,
    columns=fullflexible,
    frame=none
}

\vspace{3cm}

\title{\textbf{SSL REPORT} \\ \Large Mini Game Hub}
\author{Majid Husain \\ Vedansh Garg}
\date{} % rm -F date
\begin{document}

\maketitle

\vspace{1cm}

\tableofcontents
\newpage

\section{Introduction}
This is the Spring (2026) semester project for the course SSL (Software and Systems Lab).

\begin{itemize}
    \item The project integrates Bash for user authentication and Python (Pygame) for GUI gameplay.
    \item Players log in and their login information is stored securely with encryption.
    \item The games in the hub are:
    \begin{enumerate}
        \item TicTacToe
        \item Othello
        \item Connect4
        \item TBD
    \end{enumerate}
    \item At the end of the game, the result is stored and a leaderboard is displayed with visual representations using Matplotlib.

    \item The players are then prompted with an option to play another game or quit.
\end{itemize}

\section{Running Instructions}
\begin{itemize}
    \item The bash script is run using:
    \begin{verbatim}
    bash main.sh
    \end{verbatim}

    \item It prompts the users to input usernames along with their passwords.

    \item If a user is not registered, the script prompts to register a new username with a password.

    \item If both players successfully log in, the script internally calls:
    \begin{verbatim}
    python3 game.py <player1> <player2>
    \end{verbatim}

    \item The game hub then starts, and the players may choose any game from the menu.
\end{itemize}

\section{Modules Used}

\subsection{NumPy}
\href{https://www.geeksforgeeks.org/numpy/python-numpy/}{\color{blue}NumPy} is a Python library used for fast and efficient numerical computations and array operations. \\
It is used to represent the game boards and perform win-condition checks.

\subsection{Matplotlib}
\href{https://www.geeksforgeeks.org/python/python-introduction-matplotlib/}{\color{blue}Matplotlib} is a Python library used for creating plots and visualizations. \\
It is used to display graphs and charts of player stats after the game ends.

\subsection{datetime}
\href{https://www.geeksforgeeks.org/python/python-datetime-module/}{\color{blue}datetime} is a Python module used for handling date and time-related operations. \\
It is used for in-game timings and adding game dates to \texttt{history.csv}.

\subsection{Pygame-CE}
\href{https://pygame-gui.readthedocs.io/en/latest/quick_start.html}{\color{blue}Pygame-CE} is a library used for building graphical user interfaces (GUI) and games. \\
It is used to implement the GUI and player interactions within the game.

\subsection{sys}
The \href{https://www.geeksforgeeks.org/python/python-sys-module/}{\color{blue}sys module} provides access to command-line arguments and system-specific parameters. \\
It is used to read player usernames passed as arguments from \texttt{main.sh} to \texttt{game.py}.



\section{Directory Structure}

\begin{lstlisting}
project-root/
    Report/
        Report.tex
        Makefile
    hub/
        main.sh
        game.py
        leaderboard.sh
        games/
            tictactoe.py
            othello.py
            connect4.py
        users.tsv
        history.csv
\end{lstlisting}



\section{File descriptions}

\subsection{bash files}

\subsubsection{main.sh}

The bash script which has to be called to start the game hub
\begin{verbatim}
    bash main.sh
\end{verbatim}
The script prompts to log in two players with their usernames and passwords while doing necessary checks. If the login is successful, the script calls \texttt{game.py} with player usernames as command line arguments.
\begin{verbatim}
    python3 game.py <player1> <player2>
\end{verbatim}

\subsubsection{leaderboard.sh}
On completion of the game, \texttt{game.py} calls \texttt{leaderboard.sh}. The script reads data from \texttt{history.csv} and prints player stats for each game as formatted tables.

\subsection{python files}
\subsubsection{game.py}
\texttt{game.py} is the main Python file which reads player names passed by \texttt{main.sh} and manages the Python-side flow. The game menu, gameplay, post-game recording, and analytics are all managed by \texttt{game.py}.
The base class \textbf{Game} is also defined in this file.


\subsubsection{Games}

\begin{itemize}
    \item \texttt{tictactoe.py}

    The class TicTacToe and the run function is defined here. \texttt{game.py} calls \texttt{tictactoe.py} if the user selects it from the GUI. The run function returns the winner.

    \item \texttt{othello.py}

    The class Othello and the run function is defined here. \texttt{game.py} calls \texttt{othello.py} if the user selects it from the GUI. The run function returns the winner.

    \item \texttt{connect4.py}

    The class Connect4 and the run function is defined here. \texttt{game.py} calls \texttt{connect4.py} if the user selects it from the GUI. The run function returns the winner.
\end{itemize}

\subsection{Data files}
    \subsubsection{users.tsv}
    A tab-separated file where each line stores the username of a registered player along with the password, stored not as plain text but hashed with SHA-256.

    \subsubsection{history.csv}
    A comma-separated file that stores the winner, loser, date, and game name as a new row after each game ends.
    This data is then read by \texttt{leaderboard.sh} to print formatted tables of player stats.

    \subsection{Report}

    \subsubsection{Report.tex}
    This \texttt{.tex} document contains the detailed report explaining the project. It is compiled into a PDF.

    \subsubsection{makefile}
    In the same folder as \texttt{Report.tex}, there is the makefile used to compile it into a PDF.
\end{document}